\documentclass[11pt]{article}

\newcommand{\cnum}{Math 245B}
\newcommand{\ced}{Winter 2026}
\newcommand{\ctitle}[4]{\title{\vspace{-0.5in}\cnum, \ced\\Week #1: #2}\author{\vspace{-0.35in}\\#3, #4}}
\usepackage{enumitem}
\usepackage[usenames,dvipsnames,svgnames,table,hyperref]{xcolor}
\usepackage{amsmath, amsfonts}
\usepackage{graphicx} % Required for inserting images
\usepackage{float}
\usepackage{listings}
\usepackage{xcolor}
\usepackage{hyperref}
\usepackage{comment}

\renewcommand*{\theenumi}{\alph{enumi}}
\renewcommand*\labelenumi{(\theenumi)}
\renewcommand*{\theenumii}{\roman{enumii}}
\renewcommand*\labelenumii{\theenumii.}

\definecolor{codegreen}{rgb}{0,0.6,0}
\definecolor{codegray}{rgb}{0.5,0.5,0.5}
\definecolor{codepurple}{rgb}{0.58,0,0.82}
\definecolor{backcolour}{rgb}{0.95,0.95,0.92}
\definecolor{header}{HTML}{205459}
\definecolor{solution}{HTML}{204d52}
\DeclareMathOperator{\spt}{spt}

\newcommand{\solution}[1]{{{\textcolor{header}{Solution:} \textcolor{solution}{#1}}}}
\setlist[enumerate]{label=(\alph*)}

\lstdefinestyle{mystyle}{
    backgroundcolor=\color{backcolour},
    commentstyle=\color{codegreen},
    keywordstyle=\color{magenta},
    numberstyle=\tiny\color{codegray},
    stringstyle=\color{codepurple},
    basicstyle=\ttfamily\footnotesize,
    breakatwhitespace=false,
    breaklines=true,
    captionpos=b,
    keepspaces=true,
    numbers=left,
    numbersep=5pt,
    showspaces=false,
    showstringspaces=false,
    showtabs=false,
    tabsize=2
}


\begin{document}
\ctitle{\#2}{}{Asher Christian}{006-150-286}
\date{}
\maketitle
\section{Reminder}
$\mu$ and $\nu$ are Radon measures
\section{Theorem [ Lebesgue's decomposition ] } There are Radon measures  $\nu_{abs}$ and  $\nu_s$ such that  $\nu = \nu_{abs} + \nu_s$.
$\nu_{abs} << \mu$ and $\nu_s \perp \mu$ and for every $A \subset \mathbb{R}^{d}$ $\mu$-measurable $\nu_{abs}[A] = \int_{A}^{}D_\mu \nu(x) \mu(dx)$.
\[
    \nu_{abs}[A] = \int_{A}^{} D_\mu \nu_{abs} d \mu
\] 
Proof: We assume without loss of generality that $\mu$ and $\nu$ are finite. Let $B$ be a Borel set minimizing the folowing problem
 \[
     \inf_{C}\{\nu[C] : C \subset \mathbb{R}^{d} \text{ Borel } \mu[C^{c}] = 0\}
\] 
Set $\nu_{abs} = \nu|_B$ and $\nu_s = \nu|_{B^{c}}$. Consider the following radon measure. Then 
\[
    \nu_{abs} << \mu
\] 
By assumption $\mu[B^{c}] = 0$ and $\mu_s[B] = \mu[B \cap B^{c} ] 0$ and so $\mu \perp \nu_s$.
We claim that $\nu_{abs} << \mu$. Proof of the claim. we want to show that $\mu[E] = 0 \implies \nu_{abs}[E] = 0$.
Let $E \subset \mathbb{R}^{d}$ be such that $\mu[E] = 0$. Let $E' \subset \mathbb{R}^{d}$ be a borel set such that $E \subset E'$ and $\mu[E] = \mu[E']$. Replacing $E$ by $E'$ if necessary,
we assume wlog that  $E$ is Borel. Let
\[
    \mathcal{E} = \{c \subset \mathbb{R}^{d} \text{ Borel } | \mu[C^{c}] = 0\}
\] 
by assumption $E^{c} \in \mathcal{E}$. But
\[
    \nu[B] = \nu[B \cap E] + \nu[B \cap E^{c}] > \nu[B \cap E^{c}] \text{ unless $\nu[B \cap E] = 0$.}
\] 
and so, $\nu[B \cap E] = 0$ otherwise we will have that  $\nu[B] > \nu[B \cap E^{c}] $ which contradicts the minimality property of $B$ since $\mu[(B\cap E^{c})^{c}] = 0$. so
\[
    0 = \nu[B \cap E] = \nu_{abs}[E]
\] 
we claim that
\[
D_\mu \nu_s = 0 \qquad \mu \;\; a.e
\] 
Proof of the claim: Since $\mu[B^{c}] = 0$ it suffices to show that the measure $\mu[B \cap \{D_\mu \nu_s > 0\}] = 0$. For each $\alpha > 0$
\[
    C_\alpha := B \cap \{D_\mu \nu_s \ge \alpha\} \subset B \implies 0 = \nu_s[C_\alpha] \ge \alpha \mu[C_\alpha] \implies \mu[C_\alpha] = 0  \;\; \forall \alpha > 0
\] 
Thus proves that the measure $\mu[\{D_\mu \nu_s > 0\} \cap B] = 0$ and concludes the proof of the claim.\\
Let $A \subset \mathbb{R}^{d}$ be $\mu $-measurable.We have that 
\[
    D_\mu\nu = D_\mu \nu_{abs} + D_\mu \nu_s = D_\mu \nu_{abs} \qquad \mu \;\; a.e
\] 
since
\[
    \nu[B_r(x)] = \nu_{abs}[B_r(x)] + \nu_s[B_r(x)]
\] 
By the Radon-Nikodym Theorem
\[
    \nu_{abs}[A] = \int_{A}^{} D_\mu \nu_{abs} d \mu = \int_{A}^{}D_\mu \nu d \mu
\] 

\section{def}
\emph{
    If $f : \mathbb{R}^{d} \rightarrow [-\infty,\infty]$ is $\mu$-measurable and $p > 0$ we say that  $f$ is $p$-locally summable with respect to  $\mu$ and write $f \in L^{p}_{loc}( \mathbb{R}^{d}, \mu)$
    if $x_kf \in L^{p}( \mathbb{R}^{d}, \mu)$ for every $k \subset \mathbb{R}^{d}$ compact
}
\\
\emph{
    Notation If $E \subset \mathbb{R}^{d}$ and $\mu[E] > 0$, we write $\int_{E}^{}f d \mu = \frac{1}{\mu[E]} \int_{E}^{}f d \mu$ where the first integral has a line through it
}
\section{Lemma}
\emph{
    Let $f \in L_{loc}^{1}( \mathbb{R}^{d}, \mu)$ then for $\mu$ almost every $x \in \mathbb{R}^{d}$ 
    \[
    f(x) = \lim_{r\to 0} \frac{1}{\mu(B_r(x))}\int_{B_r(x)}^{}f d\mu
    \] 
}
Proof. We prove the result for $f^{\pm}$ where $f^{\pm} = \max\{0,\pm f\}$. This means we can assume $f \ge 0$. Set
 \[
     \nu[A] = \inf_B \{\int_{B}^{}f d \mu : B \subset \mathbb{R}^{d} \text{ is Borel and } A \subset B\}
\] 
Note that $\nu$ is a Radon measure and $\nu << \mu$. Furthermore $\nu[B] = \int_{B}^{}f d \mu$ if $B \subset \mathbb{R}^{d}$ is borel.
By the Radon Nikodym Theorem
\[
    \nu[B] = \int_{B}^{}D_\mu \nu d \mu \qquad \text{if $B \subset \mathbb{R}^{d}$ is Borel}
\] 
by previous we get
\[
D_\mu \nu = f \qquad \mu \;\; a.e
\] 
if
\[
\int_{A}^{}f = 0 \qquad \forall A \implies f = 0
\] 
Thus $D_\mu \nu(x) = \lim_{r\to 0} \frac{\nu[B_r(x)]}{\mu[B_r(x)]} = \lim_{r\to 0} \frac{1}{\mu[B_r(x)]} \int_{B_r(x)}^{}f d \mu$

\section{Friday}
Reminder  $\mu$ and $\nu$ are Radon measures. If  $f \in L^{1}_{loc}( \mathbb{R}^{d}, \mu)$ then $\lim_{r\to 0^{+}} \frac{1}{B_r(x)} \int_{B_r(x)}^{}f d \mu = f(x)$ $\mu$ almost everywhere.
\\
Corollary:
\emph{
    Let $p \in [1, +\infty)$ and $f \in (L^{p}( \mathbb{R}^{d}, \mu))$ Then
    \[
        (*) \qquad \lim_{r\to 0^{+}} \frac{1}{\mu(B_r(x))}\int_{B_r(x)}^{}|f-f(x)|^{p}d \mu = 0 \qquad \mu \text{a.e}
    \] 
}
Set $g_x(y) = |f(y)-f(x)|^{p}$ $g_x \in L^{1}_{loc}$ and the theorem holds for $y \in R^{d} \setminus N_x$ where $\mu[N_x] = 0$.\\
Proof: Let $(r_i)_{i=1}^{\infty}$ be a dense subset of $ \mathbb{R}$. Set $A_i = \{x \in \mathbb{R}^{d}: \lim_{r\to 0^{+}} \frac{1}{\mu[B_r(x)]}\int_{B_r(x)}^{}|f-r_i|^{p} d \mu = |f(x)-r_i|^{p}\}$ and let
\[
    N_i = \mathbb{R}^{d} \setminus A_i \qquad N = \bigcup_{i=1}^{\infty}N_i
\] 
Since $\mu[N_i] = 0$ then $\mu[N] = 0$. Let $x \in \mathbb{R}^{d} \setminus N$ we want to show that for every $\epsilon > 0$
\[
    \limsup_{r\to 0^{+}} \frac{1}{\mu[B_r(x)]}\int_{B_r(x)}^{}|f-f(x)|^{p}\mu(dy) \le \epsilon
\] 
Choose $i$ such that  $|f(x)-r_i| < 2^{-p}\epsilon$. We have $|f-f(x)|^{p} \le 2^{p-1}(|f-r_i|^{p}+|r_i - f(x)|^{p}) \le 2^{p-1}(|f-r_i|^{p}+ \epsilon) \le 2^{p-1}|f(x)-r_1|^{p} + \frac{\epsilon}{2}$ 
we have if $p \ge 1$ then $z \rightarrow |z|^{p}$ is convex and so 
\[
|a+b|^{p} \le c_p (|a|^{p} + |b|^{p})
\] 
We have that 
\[
\frac{1}{\mu(B_r(x))}\int_{B_r(x)}^{}|f-f(x)|^{p}d \mu \le 2^{p-1}\frac{1}{\mu(B_r(x))}\int_{B_r(x)}^{} |f(x)- r_1|^{p} d \mu + \frac{\epsilon}{2}
\] 
since $x \not\in N$ this implies that 
 \[
\limsup_{r\to 0^{+}}\frac{1}{\mu(B_r(x))}\int_{B_r(x)}^{}|f-f(x)|^{p} d \mu \le \frac{\epsilon}{2}
\] 
\\
\emph{
    Definition: when $(*)$ holds and  $p=1$ we call  $x$ a Lebesgue point of $f$ with respect to $\mu$.
}

\section{Exercise}
\emph{
    Let $E \subset \mathbb{R}^{d}$ be a $\mu$-measurable set. Then
    \[
        \lim_{r\to 0^{+}} \frac{\mu[E \cap B_r(x)] }{\mu[B_r(x)]} = 
        \begin{cases}
            1 & \mu \;\; a.e. \text{ on $E$}\\
            0 & \mu \;\; a.e. \text{ on $E^{c}$ }
        \end{cases}
    \] 
}
Proof: Let $f = 1_E$ then
 \[
     \frac{1}{\mu(B_r(x))} \int_{ B_r(x)}^{}|f-f(x)|^{p} d \mu =  \frac{1}{\mu[B_r(x)]}
     \begin{cases}
         \int_{B_r(x)}^{}(1-f) d \mu = \mu[B_r(x)] - \mu[E \cap B_r(x)] & x \in E\\
         \int_{B_r(x)}^{}f d \mu = \mu[B_r(x) \cap E] & x \not\in E
     \end{cases}
\] 
\\
Notation:
\begin{enumerate}
    \item If $f: \mathbb{R}^{d} \rightarrow \mathbb{R}^{d}$ then $\spt(f) = \overline{\{x \in \mathbb{R}^{d}: f(x) \ne 0\}}$
    \item $C_c( \mathbb{R}^{d}, \mathbb{R}^{d})  = \{f \in C( \mathbb{R}^{d}, \mathbb{R}^{d}) : \spt(f) \text{ is compact}\}$
\end{enumerate}
Let $p \in [1, \infty]$ and set $p' = \frac{p}{p-1}$ and $\frac{1}{p} + \frac{1}{p'} = 1$. The Riesz Representation says
If you take $(L^{p}( \mathbb{R}^{d}. \mu))' = L^{p'}( \mathbb{R}^{d}, \mu)$ if $p < +\infty$. Rephrased if 
\[
    L : L^{p}( \mathbb{R}^{d}, \mu) \rightarrow \mathbb{R} \qquad \text{ is linear and continuous}
\] 
Then $L$ can be represented via an element $g \in L^{p'}( \mathbb{R}^{d}, \mu)$
\[
L(f) = \int_{ \mathbb{R}^{d}}^{}f g d \mu = \langle f, g \rangle_\mu
\] 
counter example when $p = \infty$ and $L : L^{\infty}( \mathbb{R}^{d}, \mu)$ $|Lf| \le c ||f||_{L^{\infty}}$ for all $f$ and we cannot find $g \in L^{1}( \mathbb{R}^{d}, \mu)$ such that $L(f) = \int_{ \mathbb{R}^{d}}^{}fg d \mu$ 
\[
L(f) = f(0)
\] 
Topology on $C_c( \mathbb{R}^{d}, \mathbb{R}^{D})$ We say that $(f_k)_k \subset C_c( \mathbb{R}^{d}, \mathbb{R}^{D}) $converges to $f \in C_c( \mathbb{R}^{d}, \mathbb{R}^{D}) $ if there exists
$R > 0$ such that 
 \[
     \spt (f_k) \subset B_R(0) \text{ and } \limsup_{k\to +\infty} ||f-f_k||_{C( \mathbb{R}^{d}, \mathbb{R}^{D})} = 0
\] 
Note that $\spt f \subset B_R(0)$ and $f \in C_c( \mathbb{R}^{d}, \mathbb{R}^{D})$
\section{Definition}
\emph{
    We say that a linear map $L : C_c( \mathbb{R}^{d}, \mathbb{R}^{D})$ is continuous if 
    \[
        ||L|_{C(B_n(0), \mathbb{R}^{d})} || < +\infty \forall n
    \] 
    for all $n$ there exists $c_n$ such that
    \[
        |Lf| \le c_n ||f||_{C( \mathbb{R}^{d}, \mathbb{R}^{d})}
    \] 
    If the suppoort of $f$ is a subset of $B_n(0)$
}
Example: Let $\sigma : \mathbb{R}^{d} \rightarrow S^{D-1} = \{x \in \mathbb{R}^{d} : ||x|| = 1\}$ be $\mu$-measurable and define
\[
    (**) \qquad L(f) = \int_{ \mathbb{R}^{d}}^{}\langle \sigma, f \rangle d \mu \qquad \langle \sigma, f \rangle = \sum_{i=1}^{D}f_i \sigma_i
\] 
Then $L$ is linear and if $\spt f \subset B_n(0)$ then 
\[
    |L(f)| \le \int_{B_n(0)}^{}||\sigma(x)||||f(x)|| d \mu \le ||f||_{C( \mathbb{R}^{d}, \mathbb{R}^{D})} \mu[B_n(0)]
\] 
our goal will be to show that $(**)$ gives the complete lsit of  $(c_c( \mathbb{R}^{d} , \mathbb{R}^{D}))^{*}$
\\
Assume $D = 1$ then
 \[
L(f) = \int_{}^{}f d \nu
\] 
\end{document}
